\documentclass[ekobook.tex]{subfiles}

\begin{document}\setWPP
\setlength{\multicolsep}{0.5em}% Remove vertical space before and after
\setlength{\columnsep}{0.5em}% Remove space between columns
\setlength{\leftskip}{0em}%
\setlist[itemize]{
  itemsep=0pt,
  parsep=0pt,      % space between paragraph lines inside one item
  topsep=0.5ex,    % space before/after the whole list (usually keep a little)
  partopsep=0pt
}
\lsection{Logistika a zásobování}
\qaw{Vysvětli pojem logistika}{%
Logistika je věda, jejímž úkolem je zajistit, aby bylo správné \textbl{zboží} ve správném \textbl{čase}, ve správném \textbl{množství}, ve správné \textbl{kvalitě} na správném \textbl{místě} a se správnými \textbl{náklady}.
}
\qaw{Jak se dělí plány}{%
\begin{wnumbered}
\wtem{dlouhodobé plány}
nad pět let (obvykle do deseti, Japonci klidně plánují i na padesát let); obsahují způsob dosažení \texthl{vize firmy}, \texthl{podíl na trhu}; velmi obecné
\wtem{střednědobé plány}
mezi jedným rokem až pěti; konkrétní plány \texthl{příjmů a výdajů}, výstavba, nákup zařízení, ...
\wtem{krátkodobé plány}
do jednoho roku; \texthl{velmi podrobné}; plánovaní výroby, nákup materiálů, činnost zaměstnanců, ...
\end{wnumbered}
}
\vspace*{1em}
\wbreak
\qaw{Co je cílem logistických strategií}{%
Cílem je dosáhnout \texthl{výhody nad konkurencí} (nižší náklady, vyšší flexibilita, lepší služby zákazníkům).
}
\qaw{Vysvětli strategie:}{%
\begin{multicols}{2}
\textbl{ABC}\\
Nakupovaný materiál se rozdělí do tří skupin.
\begin{enumerate}[label={\color{WPP}\textbf{\Alph*)}}]
\wtem{materiál, který používáme nejčastěji}
tvoří 60--80 \% spotřeby (u pekárny mouka, droždí, cukr, mléko)
\wtem{položky méně používané}
tvoří asi 15--20 \% spotřeby (lehce nahraditelný, méně používaný; kmín, sezam)
\wtem{druhy, které jsou snadno získatelné, či se je nevyplatí skladovat}
kancelářské potřeby
\end{enumerate}
\vspace*{1em}
\textbl{JUST IN TIME}\\
Nakupujeme materiál \texthl{dle aktuální poptávky}, materiál dáme \texthl{rovnou do procesu výroby}, když není poptávka, nenakupujeme materiál, tedy \texthl{nic neskladujeme}.
\vspace*{1em}\\
\textbl{MRP}\\
(Material Requirements Planning) Postupujeme podle \texthl{přesných plánů}.
Známe přesný sled postupu výroby a materiál nakupujeme vždy před zahájením dané etapy;
\texthl{skladujeme pouze minimum}.
\columnbreak\\
\textbl{KANBAN}\\
Materiál je označen kartičkami (s čipy); nesmí se spotřebovat víc materiálu, než je na kartičce; \texthl{systém snižuje náklady na správu, lidské zdroje a čas}; slouží k signalizaci stavu zásob a rozpracované výroby.
\vspace*{1em}\\
\textbl{JIDOKA}\\
Výrobní proces je \texthl{automatizován stroji}, které se zastaví při poruše, zaměstnanci jsou pověřeni pouze běžnou kontrolou výroby, \texthl{poruchy musí opravovat specializovaní pracovníci}; zajišťuje to bezpečnost, kvalitu a efektivnost práce.
\end{multicols}
}
\vspace*{0.5em}
\wbreak
\vspace*{1em}
\begin{multicols}{2}
\qaw{Vysvětli:}{%
\texthl{minimální \wx{ } maximální zásoba} -- zásoba těsně před dodávkou \wx{ } po dodávce materiálu\\
\texthl{optimální zásoba} -- zajišťuje chod výroby bez přerušení\\
\texthl{pojistná zásoba} -- zásoba pro případ výpadku dodávek (třeba vánice pro byznysy na horách)\\
\texthl{technická zásoba} -- zásoba, co musí teprve dozrát (dřevo musí vyschnout, sýr dozrát)
}
\qaw{Jak dlouho je oběžný majetek v podniku}{%
Zpravidla \texthl{do jednoho roku} a v průběhu používání \texthl{mění svou podstatu} (peníze → materiál → výrobky → pohledávky → peníze).
}
\qaw{Popiš koloběh majetku}{%
Za \texthl{finanční majetek} nakoupíme materiál či zboží → materiál zpracujeme do \texthl{výrobků} → prodáme odběratelům → vznikne \texthl{pohledávka} → po úhradě máme \texthl{více finančního majetku} než na začátku.
}
\qaw{Charakterizuj zásoby ve výrobě}{%
\begin{enumerate}[label=\textcolor{WPP}{\textbf{\alph*)}}]
\item \textbl{materiál, pomocné látky, provozovací látky, obaly} -- vstupy do výroby
\item \textbl{polotovary} -- částečně dokončené výrobky určené k dalšímu zpracování nebo prodeji
\item \textbl{nedokončená výroba} -- rozpracované výrobky, které ještě nejsou hotové
\item \textbl{výrobky} -- \texthl{finální produkty vlastní výroby} určené k prodeji
\item \textbl{zboží} -- \texthl{odkoupené výrobky} od dodavatelů určené k přeprodeji
\item \textbl{náhradní díly} -- součástky pro opravy strojů a zařízení
\item \textbl{obaly} -- materiál sloužící k balení a ochraně zboží
\end{enumerate}
}
\qaw{Vysvětli rozdíl mezi výrobkem a zbožím}{%
\textbl{výrobek} -- \texthl{finální produkt vlastní výroby} určený k prodeji\\
\textbl{zboží} -- \texthl{odkoupený výrobek} od dodavatele, který \texthl{přeprodáváme s přirážkou}
}
\qaw{V jaké podobě můžou být peníze}{%
\begin{itemize}[label={\color{WPP}\faSlackHash}]
\item \textbl{hotové peníze} -- \texthl{pokladna}, bankovky a mince
\item \textbl{bezhotovostní peníze} -- \texthl{peníze na bankovním účtu}
\item \textbl{ceniny} -- poštovní známky, kolky, stravenky, telefonní karty
\end{itemize}
}
\qaw{Co to jsou pohledávky}{%
\texthl{Právo na budoucí příjem} peněz (faktury za dodané zboží/výrobky, které ještě nebyly uhrazeny). Vznikají po prodeji odběrateli.
}
\qaw{Přiřaď:}{%
materiál \texthl{→} látka na šaty\\
pomocné látky \texthl{→} nitě, knoflíky, zip\\
provozovací látky \texthl{→} olej na mazání stroje\\
obaly \texthl{→} lahev na víno\\
náhradní díly \texthl{→} náhradní součástka ke stroji
}
\end{multicols}
\wpage
\end{document}